\section{Related Work}
\label{sec:related}

Our work sits at the intersection of three active areas: consistency-based uncertainty quantification, token log-probability signals for hallucination detection, and LM-as-a-judge evaluation methodology. We review each in turn, highlighting how our contribution bridges gaps that prior work leaves open.

\subsection{Consistency-Based Uncertainty Quantification}

Self-consistency as a reliability signal was established by \citet{wang2022self}, who showed that sampling $N$ outputs and selecting the most frequent answer reliably improves reasoning task performance --- demonstrating that output diversity correlates with uncertainty. \citet{manakul2023selfcheck} operationalized this as SelfCheckGPT: a black-box hallucination detection method that scores consistency across $N$ samples using NLI-based contradiction detection, BERTScore, or n-gram overlap. The NLI variant achieves NonFact AUC-PR of 92.50 vs.\ 83.21 for log-probability signals on WikiBio, establishing that consistency-based signals outperform log-prob signals in document generation tasks.

The foundational formalization of consistency-based UQ is Semantic Entropy (SE) \citep{kuhn2023semantic}, which defines uncertainty as the entropy over semantic equivalence classes:
\[
H(p(C|x)) = -\sum_{C} p(C|x) \log p(C|x),
\]
where $p(C|x) = \sum_{s \in C} p(s|x)$ sums token probabilities of all outputs in the same NLI-defined cluster. \citet{kuhn2023semantic} demonstrate that SE outperforms raw token entropy on TriviaQA and BioASQ. Our work directly builds on this formulation ($N=5$, bidirectional DeBERTa-MNLI NLI clustering per Kuhn et al.\ Eq.~3).

More recent work extends SE to address its limitations in longer-form generation. \citet{nguyen2025beyond} (SNNE) introduces pairwise SE computation to handle multi-sentence answers where the original SE formulation weakens. This directly motivates our scope restriction to short-answer factual QA (TriviaQA, $\leq$50 token answers), where original SE remains the strongest consistency-based signal. \citet{bouchard2025uncertainty} (UQLM) benchmark 50+ UQ methods across 24 model--dataset combinations and find that NLI-based black-box scorers are best in 13/24 AUROC scenarios --- providing strong external validation that SE-family signals are competitive.

\textbf{Gap:} None of these works directly measure the Pearson correlation between SE and min\_logprob on open-weight LLMs. They report standalone AUROC comparisons, not cross-signal statistical independence. Our work closes this gap.

\subsection{Token Log-Probability Signals for Hallucination Detection}

Token log-probabilities are the most computationally efficient UQ signal class, requiring only a single forward pass. \citet{malinin2021uncertainty} formalize sequence-level uncertainty measures from log-probabilities. Min\_logprob has been empirically shown to achieve strong single-pass baseline performance on TriviaQA dev with Llama-3.1-8B (our prior internal pipeline experiment h-m1).

The relationship between first-token confidence and SE was characterized by \citet{gabriel2026first}, who finds $r = 0.54$--$0.76$ between first-token probability and semantic agreement in Llama-3.1-8B, TriviaQA --- suggesting that for first-token signals, moderate correlation with SE exists. A critical distinction: min\_logprob is NOT first-token confidence. It is the minimum over ALL tokens in greedy decode, capturing arbitrarily late positions. This operational difference drives near-zero Pearson $r$ in our measurements ($|r| = 0.026$).

\citet{dey2025uncertainty} (UAF Ensemble) show that fusing UQ signals improves factual accuracy by 8\%, further motivating the ensemble direction.

\textbf{Gap:} No prior work directly measures the statistical independence of min\_logprob and SE on open-weight LLMs under a controlled evaluation protocol. Our $|r|=0.026$ measurement fills this gap.

\subsection{LM-as-a-Judge Evaluation Methodology}

\citet{santilli2025revisiting} (ACL 2025) show that AUROC rankings for hallucination detection methods are systematically distorted by response length bias --- ROUGE-L evaluation is most severely affected. They recommend LM-as-a-judge with cross-model design as the most human-aligned correctness function. \citet{xiong2023llms} benchmark consistency-based signals as the most reliable black-box UQ method overall.

We implement the cross-model judge recommendation directly: Qwen-2.5-7B-Instruct evaluates Llama-3.1-8B-Instruct outputs on TriviaQA dev. Our Spearman $\rho(\text{SE}, \text{judge}) = -0.026$ directly validates this design choice: the circularity threshold of $|\rho| < 0.40$ is satisfied with substantial margin.

\textbf{Our Position:} We complement prior work by characterizing the cross-signal statistical relationship (SE vs.\ min\_logprob independence) that justifies their combination in an ensemble, using the evaluation best practices that prior work has established.
